\section{Реализация, интеграция и тестирование}\label{sec:implementation}

\subsection{Реализация даунтайм-этапа}\label{subsec:impl-mvp}

\subsubsection{Оркестратор внутри service-staas}\label{subsubsec:impl-orchestrator}

Оркестратор реализован как фоновая горутина внутри \textit{service-staas}, запускаемая параллельно с HTTP-частью сервиса.
Воркер периодически опрашивает БД платформы и забирает в обработку бакеты в активных миграционных статусах (\code{MIGRATION\_SCHEDULED}, \code{MIGRATION\_IN\_PROGRESS}, \code{MIGRATION\_VERIFYING}).

Для каждого взятого бакета воркер выполняет действия текущей фазы (создание целевого бакета, запуск копирования, сверка, переключение DNS) и по завершении фазы переводит бакет в следующий статус.
Транзиентные ошибки (медленный кластер, временная недоступность Consul или PostgreSQL) не переводят бакет в терминальный статус: воркер откладывает следующую попытку по фиксированному расписанию (1, 2, 5, 10 минут, далее 10-минутный интервал) до успешного завершения фазы.
\code{MIGRATION\_FAILED} выставляется только при неуспешной сверке на фазе \code{MIGRATION\_VERIFYING}~--- это означает, что данные перенесены с расхождениями и требуется ручной разбор.

\subsubsection{Интеграция с rclone}\label{subsubsec:impl-rclone}

Рядом с \textit{service-staas} поднимается процесс \texttt{rclone rcd}~--- режим Remote Control daemon, в котором \gls{rclone} принимает HTTP-команды на запуск, опрос и отмену задач~\cite{rclone-rc}.
Воркер стартует копирование вызовом \texttt{sync/sync} с указанием исходного и целевого \gls{s3}-бэкендов, получает идентификатор задачи, далее периодически опрашивает её через \texttt{job/status} и при необходимости отменяет вызовом \texttt{job/stop}.

Параметры подключения к каждому из бэкендов задаются в конфигурационном файле rclone в виде отдельных секций, именованных по имени кластера; учётные данные подставляются воркером из \gls{vault} перед стартом задачи в шаблон конфигурации.

Пример формирования шаблона конфигурации rclone в воркере приведён в листинге~\ref{lst:rclone-conf}.

\begin{listing}[ht]
\begin{minted}{go}
const rcloneConfigTemplate = `
[{{ .SrcCluster }}]
type = s3
provider = Ceph
endpoint = {{ .SrcEndpoint }}
access_key_id = {{ .SrcAccessKeyID }}
secret_access_key = {{ .SrcSecretAccessKey }}

[{{ .DstCluster }}]
type = s3
provider = AWS
endpoint = {{ .DstEndpoint }}
access_key_id = {{ .DstAccessKeyID }}
secret_access_key = {{ .DstSecretAccessKey }}
`
\end{minted}
\caption{Шаблон конфигурации rclone, формируемый воркером}
\label{lst:rclone-conf}
\end{listing}

\subsection{Реализация бесшовного этапа}\label{subsec:impl-seamless}

\subsubsection{Прокси s3-shifter}\label{subsubsec:impl-shifter}

\textit{s3-shifter} реализован как HTTP-сервер на Go, принимающий запросы по \gls{s3}-протоколу.
Входящий запрос разбирается до уровня, достаточного для маршрутизации~--- определяются тип операции, имя бакета и ключ объекта; тело запроса при этом не вычитывается в память, а потоково пересылается дальше через \texttt{io.Copy}.

Список бакетов, находящихся под миграцией, и адреса исходного и целевого кластеров \textit{s3-shifter} периодически забирает из \textit{service-staas}.

Перед каждой записывающей операцией на объект под миграцией берётся блокировка (подраздел~\ref{subsubsec:consul-locks}); по завершении операции блокировка снимается.
Чтение и листинг работают без блокировок: их корректность обеспечивается логикой маршрутизации (подраздел~\ref{subsubsec:s3-shifter}).

\paragraph{Подпись S3-запросов.}
\textit{s3-shifter} верифицирует AWS SigV4-подпись клиента и пересигнирует исходящий запрос к backend'у собственной сервисной учётной записью~--- это даёт единую трастовую границу и не зависит от того, как backend интерпретирует Host-заголовок.

\subsubsection{Воркеры Lister и Copier}\label{subsubsec:workers}

Lister и Copier реализованы как горутины внутри единого процесса мигратора, написанного на Go.
Один процесс держит один Lister и пул Copier-горутин фиксированного размера; число параллельных копирований задаётся параметром конфигурации, остальные ключи ждут своей очереди в БД.
По сути это и есть текущий механизм rate limiting: он не перегружает исходный и целевой кластеры на старте миграции, но не подстраивается под нагрузку~--- динамическое управление параллелизмом остаётся открытым направлением для развития.

Процесс мигратора разворачивается на нескольких узлах с многогигабитной сетью; общая БД и блокировки в Consul делают такое горизонтальное масштабирование безопасным.

\subsubsection{Персистентное хранилище прогресса миграции}\label{subsubsec:persistent-store}

Прогресс миграции хранится в PostgreSQL в двух таблицах: \texttt{object\_keys} с записью на каждый ключ и \texttt{migration\_state} с агрегатным статусом миграции бакета.

\paragraph{Структура \texttt{object\_keys}.}
\begin{itemize}
    \item \texttt{bucket}, \texttt{object\_key}~--- идентификатор объекта;
    \item \texttt{status}~--- состояние ключа: \code{INIT}, \code{IN\_PROGRESS}, \code{COPIED}, \code{EXIST} или \code{DELETED};
    \item \texttt{owner\_id}, \texttt{lease\_until}, \texttt{fence\_token}~--- механика аренды и защиты от устаревшего владельца;
    \item \texttt{attempt}, \texttt{updated\_at}~--- счётчик попыток и время последнего обновления.
\end{itemize}
Таблица партицируется по \texttt{bucket}, что позволяет одним \texttt{DROP PARTITION} быстро отбрасывать данные завершённой миграции.

\paragraph{Структура \texttt{migration\_state}.}
\begin{itemize}
    \item \texttt{bucket}~--- идентификатор бакета;
    \item \texttt{status}~--- технический статус: \code{IN\_PROGRESS}, \code{FINISHED}, \code{FAILED};
    \item \texttt{updated\_at}, \texttt{error\_text}~--- время последнего изменения и сообщение об ошибке.
\end{itemize}
Технический статус из этой таблицы мигратор отдаёт в \textit{service-staas}, где он отражается в полной статусной модели платформы (\code{MIGRATION\_*}).

\paragraph{Горячий путь Copier.}
Работа с \texttt{object\_keys} построена на трёх коротких SQL-операциях:

\begin{itemize}
    \item \textbf{claim}~--- выбор очередного \code{INIT}-ключа через \texttt{SELECT \ldots FOR UPDATE SKIP LOCKED}, перевод в \code{IN\_PROGRESS}, фиксация владельца, дедлайна аренды и инкремент \texttt{fence\_token};
    \item \textbf{finalize}~--- перевод ключа в терминальный статус. Условие \texttt{fence\_token = \$4} гарантирует, что обновление выполняет только тот воркер, который начал обработку: устаревший владелец просто не найдёт строки;
    \item \textbf{reclaim}~--- фоновое возвращение в \code{INIT} ключей с истёкшей арендой, после чего их подхватывает другой воркер.
\end{itemize}

Такая схема сериализует распределение работы между параллельными Copier-горутинами без отдельного брокера очередей и не удерживает блокировки PostgreSQL на время сетевого копирования (claim и finalize выполняются отдельными короткими транзакциями).
Соответствующие SQL-запросы приведены в листинге~\ref{lst:sql-hot-path}.

\begin{listing}[ht]
\begin{minted}{sql}
-- Claim объекта в работу
WITH claimed AS (
    SELECT bucket, object_key
    FROM object_keys
    WHERE bucket = $1 AND status = 'INIT'
    ORDER BY updated_at
    LIMIT 1
    FOR UPDATE SKIP LOCKED
)
UPDATE object_keys ok
SET status = 'IN_PROGRESS',
    owner_id = $2,
    lease_until = NOW() + INTERVAL '15 minutes',
    fence_token = fence_token + 1,
    attempt = attempt + 1,
    updated_at = NOW()
FROM claimed c
WHERE ok.bucket = c.bucket AND ok.object_key = c.object_key
RETURNING ok.bucket, ok.object_key, ok.fence_token;

-- Финализация только текущим владельцем поколения
UPDATE object_keys
SET status = $3, owner_id = NULL, lease_until = NULL, updated_at = NOW()
WHERE bucket = $1 AND object_key = $2
  AND status = 'IN_PROGRESS' AND fence_token = $4;

-- Reclaim просроченных задач
UPDATE object_keys
SET status = 'INIT', owner_id = NULL, lease_until = NULL, updated_at = NOW()
WHERE status = 'IN_PROGRESS' AND lease_until < NOW();
\end{minted}
\caption{SQL-операции claim/finalize/reclaim в горячем пути Copier}
\label{lst:sql-hot-path}
\end{listing}

\subsubsection{Распределённые блокировки на объектах через Consul}\label{subsubsec:consul-locks}

Блокировка реализована как Consul-сессия с конечным TTL; в текущей конфигурации используется \texttt{TTL=15m}.
Процесс, удерживающий блокировку, продлевает сессию каждые \texttt{TTL/3} (раз в 5 минут).
Если \textit{s3-shifter} или Copier упал либо потерял связь с Consul, продление прекращается, сессия истекает, и блокировка освобождается автоматически~--- объект не остаётся заблокированным навсегда.
После истечения сессии блокировка берётся заново обычной попыткой acquire; обработчик не продолжает критическую секцию со старой (просроченной) сессией.
Последовательность обработки объекта в штатном и аварийном сценариях показана на рис.~\ref{fig:lock-sequence}.

\begin{figure}[H]
    \centering
    \includegraphics[width=\textwidth]{lock-sequence}
    \caption{Sequence-диаграмма lease/fencing и Consul-lock при копировании объекта}
    \label{fig:lock-sequence}
\end{figure}

\subsection{Обоснование выбора инфраструктуры}\label{subsec:tech-choice}

Ключевые инфраструктурные компоненты системы~--- хранилище прогресса, координация блокировок и язык реализации~--- допускают разные варианты выбора.
Принятые решения опираются прежде всего на то, что уже используется в платформе \gls{staas}: дополнительный класс зависимостей повышает стоимость эксплуатации и зону отказа.

\paragraph{Хранилище прогресса~--- PostgreSQL.}
Альтернативы~--- etcd, Redis, специализированные брокеры очередей.
etcd ориентирован на конфигурацию и плохо масштабируется на десятки миллионов строк \texttt{object\_keys}, которые возникают при петабайтных бакетах.
Redis~--- in-memory, требует тщательной настройки персистентности (AOF/RDB) для критических данных миграции и не даёт нативной queue-семантики уровня \texttt{SELECT \ldots FOR UPDATE SKIP LOCKED}.
Брокер очередей (Kafka, RabbitMQ) дал бы потоковую модель «выдать ключ воркеру», но потребовал бы строить отдельный слой состояния под ретраи, fencing и партицирование~--- по сути дублировать функции БД.

PostgreSQL спокойно держит соответствующую нагрузку, имеет встроенную queue-семантику через \texttt{SELECT \ldots FOR UPDATE SKIP LOCKED}, поддерживает партицирование по бакету и уже стабильно работает в платформе.
Это снимает необходимость в новой эксплуатационной зоне ответственности.

\paragraph{Координация блокировок~--- Consul.}
Альтернативы~--- ZooKeeper, etcd, Redis Redlock.
Все три технически решают задачу распределённых блокировок, но Consul уже развёрнут в платформе и используется как DNS-сервис.
Заводить отдельный координационный сервис только ради блокировок означало бы расширить зону отказа и увеличить операционную нагрузку без новой ценности.
Сессии Consul с TTL и lock API напрямую покрывают требуемую модель блокировки на объект с автоматическим освобождением; защита от классической проблемы устаревшего владельца~\cite{kleppmann-locks} реализуется fencing token'ом на стороне БД мигратора (подраздел~\ref{subsubsec:persistent-store}).

\paragraph{Язык реализации~--- Go.}
Альтернативы~--- Python, Java, Rust.
Python неудобен для потоковой пересылки гигабайтных тел запросов через \textit{s3-shifter} и ограничен GIL при пуле Copier'ов.
Java и Rust подходят по производительности, но в платформе доминирует Go: общие SDK для Consul, Vault, PostgreSQL и S3, общая практика сборки и деплоя, общий пул разработчиков.
Go даёт нативную горутина-модель, которая прямо соответствует пулу Copier'ов и потоковой обработке прокси через \texttt{io.Copy}, и при этом снимает необходимость в отдельной экосистеме библиотек.

\subsection{Эксплуатационные свойства}\label{subsec:cross-cutting-impl}

Большая часть сквозных свойств из подраздела~\ref{subsec:cross-cutting} конкретизируется внутри соответствующих подразделов реализации (идемпотентность и восстановление~--- через механику аренды и fencing в~\ref{subsubsec:persistent-store} и~\ref{subsubsec:consul-locks}, ограничение нагрузки~--- через пул Copier в~\ref{subsubsec:workers}).
Здесь раскрываются два свойства, не укладывающиеся в рамки отдельного компонента~--- отказоустойчивость и наблюдаемость.

\textbf{Отказоустойчивость.}
Мигратор и \textit{s3-shifter} развёрнуты в нескольких дата-центрах, Consul доступен в каждом из них~--- падение отдельного ДЦ не останавливает миграцию.
Прогресс хранится во внешней персистентной БД, поэтому перезапуск любого экземпляра подхватывает миграцию с того же места.

\textbf{Наблюдаемость.}
Платформа собирает три класса показателей:
\begin{itemize}
    \item \textit{бизнес-метрики миграции}~--- доля скопированных ключей, скорость копирования, число ошибок на объект; выводятся в дашборд и позволяют отслеживать прогресс конкретной миграции;
    \item \textit{метрики \gls{slo}}~--- латентность операций через \textit{s3-shifter} с разбивкой по типу запроса, чтобы своевременно увидеть нарушение целевых порогов из подраздела~\ref{subsec:nonfunc-req};
    \item \textit{логи и распределённые трассировки}~--- для разбора конкретных инцидентов на уровне отдельного объекта или запроса.
\end{itemize}

\subsection{Тестирование}\label{subsec:testing}

Тестирование делится на три части: юнит-тесты, локальный плейграунд для проверки S3-совместимости и нагрузочные тесты \textit{s3-shifter}.

\subsubsection{Юнит-тесты}\label{subsubsec:unit-tests}

Юнит-тестами покрыты:
\begin{itemize}
    \item переходы статусов объекта (claim/finalize/reclaim) и проверка fencing token при финализации;
    \item логика маршрутизации операций в \textit{s3-shifter} для бакетов под миграцией и вне её;
    \item объединение листингов с приоритетом целевого бакета.
\end{itemize}

\subsubsection{Локальный плейграунд: совместимость с S3-протоколом}\label{subsubsec:playground}

Плейграунд развёрнут в Docker: два Ceph-кластера, \textit{service-staas}, PostgreSQL и Consul.
Клиент работает через стандартный AWS SDK; прогоняются два режима:
\begin{itemize}
    \item \textbf{обычный режим}~--- бакет не под миграцией; ответы прокси сверяются с прямыми ответами кластера на семантическое совпадение;
    \item \textbf{режим миграции}~--- бакет проводится по полному циклу \code{MIGRATION\_*}; во время копирования клиент подаёт смешанный трафик, по завершении бакеты сверяются.
\end{itemize}

Сверка использует повторный \texttt{ListObjects} на обоих бакетах и проверяет два условия:
\begin{itemize}
    \item множества ключей в source и target совпадают;
    \item для каждого ключа дата \texttt{LastModified} в target не старше, чем в source, с учётом допустимого рассинхрона часов.
\end{itemize}
Второе условие корректно благодаря правилам маршрутизации \textit{s3-shifter}: PUT идёт только в целевой бакет, DELETE~--- в оба. Значит target никогда не отстаёт от source: он содержит либо скопированную, либо более свежую клиентскую версию (листинг~\ref{lst:integrity}).

\begin{listing}[ht]
\begin{minted}{go}
func verifyByListing(srcClient, dstClient S3Client, bucket string) bool {
    const allowedClockSkew = 2 * time.Second

    srcObjects := srcClient.ListObjects(bucket) // map[key]ObjectMeta
    dstObjects := dstClient.ListObjects(bucket)

    if len(srcObjects) != len(dstObjects) {
        return false
    }

    for key, srcMeta := range srcObjects {
        dstMeta, ok := dstObjects[key]
        if !ok {
            return false
        }
        // target не должен быть старше source: PUT идёт только в target,
        // DELETE - в оба, поэтому target всегда не старше.
        if dstMeta.LastModified.Before(
            srcMeta.LastModified.Add(-allowedClockSkew),
        ) {
            return false
        }
    }

    return true
}
\end{minted}
\caption{Сверка листингов и дат модификации между источником и целью}
\label{lst:integrity}
\end{listing}

\subsubsection{Нагрузочное тестирование s3-shifter}\label{subsubsec:load-test}

Нагрузочные тесты оценивают накладные расходы прокси на горячем пути.
Генератор подаёт смешанный профиль трафика (преобладание GET, значимая доля PUT, периодические DELETE и LIST) при разных уровнях фоновой копировальной нагрузки от мигратора~--- от простаивающего до насыщающего канал.
Эталонные значения снимаются на той же конфигурации без \textit{s3-shifter} в пути; разница и есть накладные расходы прокси.

Сводка результатов приведена в таблице~\ref{tab:proxy-load-results}.

\begin{table}[H]
\centering
\footnotesize
\renewcommand{\arraystretch}{1.2}
\begin{tabular}{|p{2.8cm}|p{5.5cm}|p{5.5cm}|}
\hline
\textbf{Этап} & \textbf{Наблюдение} & \textbf{Вывод} \\
\hline
Базовый стенд до расширения ресурсов &
Кластер без прокси выдерживал GET до 15\,000 RPS и PUT до 12\,000 RPS; при включении \textit{s3-shifter} потолок упирался в ресурсы стенда. &
Узким местом была инфраструктура стенда, а не логика маршрутизации. \\
\hline
Влияние проксирования на задержку (до апгрейда стенда) &
При работе на пределе ресурсов типовая задержка росла с $\sim$40\,мс до $\sim$60\,мс. &
Оценка накладных расходов в режиме насыщения завышена. \\
\hline
Стенд после обновления железа (3 ноды) &
С прокси достигнуты GET и PUT до 15\,000 RPS, добавленная задержка 0--20\,мс. Потребление ресурсов: GET до $\sim$85~\%, PUT до $\sim$70~\% от limits; память стабильна. &
Прокси выполняет целевые требования по пропускной способности с умеренными накладными расходами по латентности. \\
\hline
\end{tabular}
\caption{Сводка результатов нагрузочного тестирования \textit{s3-shifter}}
\label{tab:proxy-load-results}
\end{table}

Тесты подтверждают, что бесшовная миграция прозрачна для клиента: ошибок клиент не получает, финальная сверка листингов и дат проходит без расхождений, а накладные расходы \textit{s3-shifter} укладываются в \gls{slo} платформы из подраздела~\ref{subsec:nonfunc-req}.

\subsection{Выводы по главе}\label{subsec:ch3-summary}

В главе подробно рассмотрена реализация спроектированного решения.
Оркестратор миграции реализован как фоновая горутина внутри \textit{service-staas}, на этапе с даунтаймом он управляет внешним процессом \gls{rclone} через rcd-API, на бесшовном этапе~--- собственным мигратором.
Мигратор разделён на роли Lister и Copier, обменивающиеся через таблицы в PostgreSQL, а согласование с прокси \textit{s3-shifter} выполняется через распределённые блокировки на объектах.

Тестирование на двухкластерном стенде Ceph подтвердило, что бесшовная миграция прозрачна для клиента: ошибок клиент не получает, финальная сверка листингов и дат сходится, накладные расходы \textit{s3-shifter} укладываются в целевые \gls{slo} платформы.
